\documentclass[../readings.tex]{subfiles}

\begin{document}

\subsection{Linear Operators and the Poisson Equation}
\label{sub:operators}

Transitioning from finite matrices to infinite-dimensional operators.

\subsubsection{Differential Operators and Discretization}

\addDef{Linear Operator}{%
Map $\mathcal{A}: X \to Y$ between function spaces satisfying:
\begin{align*}
    \mathcal{A}(\alpha u + v) = \alpha \mathcal{A}u + \mathcal{A}v.
\end{align*}
Matrices are finite-dimensional operators; $\frac{d}{dx}$ and $\nabla^2$ are infinite-dimensional.
}

\addDef{Finite Difference Stencil}{%
Approximating $u''(x)$ on a grid with spacing $h$:
\begin{align*}
    u''(x_k) \approx \frac{u(x_{k-1}) - 2u(x_k) + u(x_{k+1})}{h^2}.
\end{align*}
\textbf{Accuracy:}\enspace Second-order $O(h^2)$ via Taylor expansion.
}

\addDef{1-D Discrete Laplacian}{%
Discretizing $-d^2/dx^2$ on $[0,1]$ with $u(0)=u(1)=0$ yields the tridiagonal system $\bT \mathbf{u} = \mathbf{f}$:
\begin{align*}
    \bT = \frac{1}{h^2} \text{tridiag}(-1, 2, -1) \in \fR^{n \times n}.
\end{align*}
}

\addThm{Spectrum Convergence}{%
Eigenvalues of $\bT$ are $\lambda_k = \frac{4}{h^2} \sin^2(\frac{k\pi}{2(n+1)})$.
As $n \to \infty$, $\lambda_k \to k^2 \pi^2$ (eigenvalues of continuous $-d^2/dx^2$).
}
\label{thm:laplacian-eigenvalues}

\addPrf{%
We verify $\bT\bv_k = \lambda_k \bv_k$ by component-wise substitution.
The $j$-th entry of $\bT\bv_k$ is:
\begin{align*}
    (\bT\bv_k)_j = \frac{1}{h^2} \left[ -\sin((j-1)kh\pi) + 2\sin(jkh\pi) - \sin((j+1)kh\pi) \right].
\end{align*}
Applying the trigonometric identity $\sin(A-B) + \sin(A+B) = 2\sin A \cos B$ with $A = jkh\pi$ and $B = kh\pi$:
\begin{align*}
    (\bT\bv_k)_j &= \frac{1}{h^2} \left[ 2\sin(jkh\pi) - 2\sin(jkh\pi) \cos(kh\pi) \right] \\
    &= \frac{2}{h^2} (1 - \cos(kh\pi)) \sin(jkh\pi).
\end{align*}
Using the half-angle identity $1 - \cos\theta = 2\sin^2(\theta/2)$:
\begin{align*}
    (\bT\bv_k)_j &= \frac{4}{h^2} \sin^2\left(\frac{kh\pi}{2}\right) \sin(jkh\pi) \\
    &= \lambda_k (\bv_k)_j.
\end{align*}
As $h = 1/(n+1) \to 0$, using the small-angle approximation $\sin \theta \approx \theta$:
\begin{align*}
    \lambda_k \approx \frac{4}{h^2} \left( \frac{k\pi h}{2} \right)^2 = k^2\pi^2.
\end{align*}
}

\subsubsection{Operator Norms and Conditioning}

\addDef{Boundedness}{%
An operator is \textbf{bounded} if $\|\mathcal{A}u\|_Y \leq C \|u\|_X$ for some $C < \infty$.
\begin{itemize}
    \item Matrices are always bounded.
    \item Differential operators are \textbf{unbounded}.
\end{itemize}
}

\addNote{%
\textbf{Numerical Consequence:}\enspace Because the continuous Laplacian is unbounded, its discretization $\bT$ becomes increasingly ill-conditioned as $h \to 0$:
\begin{align*}
    \kappa(\bT) = \frac{\lambda_{\max}}{\lambda_{\min}} \approx \frac{4/h^2}{\pi^2} = O(h^{-2}).
\end{align*}
Refining the grid $2\times$ quadruples the condition number.
}

\subsubsection{Exercises}

\addExcr{%
\begin{enumerate}
    \item Build $\bT$ for $n=50$ and plot its eigenvalues vs.\ $k^2 \pi^2$.
    \item Show that $\kappa(\bT) \to \infty$ as $n \to \infty$.
    \item Implement a 2-D Poisson solver on a unit square. Verify $O(N^{3/2})$ cost for direct sparse solve.
    \item Prove the differentiation operator is unbounded on $L^2[0,1]$ using $u_k(x) = \sin(k\pi x)$.
\end{enumerate}
}

\end{document}
